% =====================================================================================================
\section{NICER and MAXI processing}\label{sec:nicermaxi}
% =====================================================================================================

% -----------------------------------------------------------------------------------------------------
\subsection{NICER/XTI: from event files to features (v9--v13)}\label{sec:nicer}
The NICER chain (\cref{fig:nicerflow}) works on cleaned event files without HEASoft. It uses the same cospectral idea
as Standard-1, with the seven NICER measurement/power units (MPUs) as detectors. HEASoft enters only for the 3C50
background (\cref{sec:3c50}). Library code: \file{v9lib.py} (parser, segments, features), \file{v10lib.py} (tail and
background proxy), \file{v11lib.py} (streaming, Lorentzians).

\begin{figure}[htbp]
\centering
% NICER/XTI pipeline (v9--v13) and the version that added each step. Copied from report_en/figures/flowcharts (styles in workflow_report.tex).
% Row 1 left to right, row 2 right to left.
\begin{tikzpicture}[x=1cm, y=1cm,
  st/.style={fc proc, text width=2.62cm, minimum height=16mm},
  badge/.style={fc ver, anchor=south west}]
  % row 1: events to per-segment cospectra
  \node[fc data, text width=2.62cm, minimum height=16mm] (ev) at (1.45,0) {cleaned event files \texttt{ni<obs>\_0mpu7\_cl}, HEASARC AWS mirror};
  \node[st] (rd)  at (4.62,0)  {streaming parser: v9 prefix ($\ge6$ segments); v10 + tail to 1\,GB; v11, v13 whole files};
  \node[st] (seg) at (7.79,0)  {good intervals (gap $>1$\,s); 128-s segments, 2--10\,keV, 1/256-s bins, $\ge2$ MPUs};
  \node[st] (bst) at (10.96,0) {segments near a type-I burst removed ($-20$\,s, $+200$\,s)};
  \node[st] (cos) at (14.13,0) {all-pairs cross-MPU cospectrum per segment};
  \foreach \a/\b in {ev/rd,rd/seg,seg/bst,bst/cos} \draw[fc arr] (\a) -- (\b);
  % row 2: observation features, quality, background
  \node[st, text width=2.62cm] (of) at (14.13,-2.4) {ratio of means: $T_1$--$T_3$, 0.1--10\,Hz state, $\nu_c$ (13 octave bins); colours $c_1,c_2$};
  \node[fc gate, text width=2.62cm, minimum height=16mm] (qg) at (10.96,-2.4) {quality: $\ge3$ segments and $\ge20$\,c\,s$^{-1}$ (2--10\,keV)};
  \node[st] (bp)  at (7.79,-2.4) {background proxy $\beta$ = 12--15 / 2--10\,keV counts; threshold 0.2687};
  \node[st] (c50) at (4.62,-2.4) {3C50 background $f$ ($<500$\,c\,s$^{-1}$): correction $1/(1-f)$, screening $f\le0.20$};
  \node[fc, draw=inkgrey, fill=boxgrey, text width=2.62cm, minimum height=16mm] (tab) at (1.45,-2.4) {feature table of the version (\file{data/v9}--\file{v13/processed})};
  \draw[fc arr] (cos) -- (of);
  \foreach \a/\b in {of/qg,qg/bp,bp/c50,c50/tab} \draw[fc arr] (\a) -- (\b);
  % side branch: Lorentzian fits
  \node[st, minimum height=11mm] (lz) at (14.13,-4.55) {one or two Lorentzians on the octave bins: $\nu_L$};
  \draw[fc arr] (of) -- (lz);
  \node[fc note, text width=7.6cm, anchor=west] at (1.45-1.38,-4.55)
     {v13 also recomputed the faint observations' features inside the 3C50 good-time intervals only (descriptive).
      $\nu_c$ is not changed by the 3C50 correction; colours and rms are.};
  % version badges
  \foreach \n/\v in {ev/v9, rd/{v9, v10, v11}, seg/v9, bst/v9, cos/v9, of/v9, qg/v9, bp/v10, c50/{v12, v13}, lz/v11c}
     \node[badge] at ([xshift=2pt, yshift=-0.4pt]\n.north west) {\v};
\end{tikzpicture}

\caption{NICER/XTI feature pipeline: top row left to right, then bottom row right to left. The v9 steps were frozen before
the observations of v11 were read, and the whole chain, including 3C50, before those of v13.}\label{fig:nicerflow}
\end{figure}

\paragraph{Step 1: read the events.} \code{v9lib.PrefixParser} turns the compressed byte stream into event arrays
(\cref{sec:get-nicer}):
\begin{enumerate}
\item \code{zlib} (gzip mode) decompresses each new piece into a buffer;
\item once the buffer holds them, the primary and \code{EVENTS} headers are parsed from 2880-byte blocks;
\item \code{row_dtype} builds a big-endian numpy record type from \code{TFIELDS}, \code{TTYPEn} and \code{TFORMn}, checking
that the row length equals \code{NAXIS1};
\item complete rows are converted and only \code{TIME}, \code{PI} and \code{DET_ID} are kept;
\item the parse is \emph{complete} when \code{NAXIS2} rows have been read.
\end{enumerate}
Events are expected in time order; if they are not, they are sorted (and flagged).

\paragraph{Step 2: good intervals and segments.} A good interval is a maximal stretch with no gap longer than 1\,s between
consecutive events of any energy (\code{v9lib.good_intervals}). Each interval $[a,b]$ is cut into consecutive 128\,s
segments starting at $a$. Only 2--10\,keV events (PI 200--999; PI is in units of 10\,eV) are used for timing. The detector
of an event is its MPU, $m=\lfloor\mathtt{DET\_ID}/10\rfloor\in\{0,\dots,6\}$. Events are binned at $1/256$\,s ($N=32768$,
Nyquist 128\,Hz). An MPU is on in a segment if each of the eight 16\,s blocks contains a 2--10\,keV event. A segment needs at
least 2 MPUs on.

\paragraph{Step 3: bursts.} The 1\,s 2--10\,keV light curve inside the good intervals (\code{v9lib.lc_1s}) is searched with
the RXTE burst rule (\cref{eq:burst}, \code{v9lib.burst_events}). Every segment that overlaps
$[t_\mathrm{start}-20\,\mathrm{s},\ t_\mathrm{end}+200\,\mathrm{s}]$ of a detected burst is skipped; the number skipped is
stored as \code{n_seg_burst_removed}.

\paragraph{Step 4: per-segment statistics (\code{v9lib.segment_stats}).} For the $n$ MPUs that are on, with mean counts per
bin $s_m$ (no background subtraction) and transforms $X_{m,j}$, the all-pairs cross power is
$\mathrm{cross}_j=\bigl(|\sum_mX_{m,j}|^2-\sum_m|X_{m,j}|^2\bigr)/(2\sinc^2(j/N))$ (\cref{eq:allpairs}). Each segment
stores a numerator $\mathrm{num}_{\mathcal B,r}=\frac{2}{N^2}\sum_{j\in\mathcal B}\mathrm{cross}_j$ for every band and for each
of the 13 octave bins, and the denominator $\mathrm{den}_r=\sum_{m<m'}s_ms_{m'}$.

\paragraph{Step 5: observation features (\code{v9lib.obs_features}).}
\begin{itemize}
\item \textbf{Timing.} Each band is combined over segments by the ratio of means with its delta-method error
(\cref{eq:ratioofmeans}), and converted to a signed square root. The bands are $T_1$, $T_2$, $T_3$, the state band
0.1--10\,Hz and the 96--128\,Hz white-noise monitor (\cref{tab:timingbands}).
\item \textbf{Colours.} \emph{All} events of all detectors inside the accepted segment windows are counted in PI 200--399
($A$, 2--4\,keV), 400--599 ($B$, 4--6\,keV) and 600--999 ($C$, 6--10\,keV); $c_1=B/A$ and $c_2=C/B$. The colours therefore
come from exactly the time windows used for timing.
\item \textbf{Count rate.} \code{rate_2_10} is the mean over segments of the 2--10\,keV counts of the MPUs that are on,
divided by 128\,s.
\item \textbf{$\nu_c$.} The 13 octave bins $k=0,\dots,12$ cover $j\in[2^k,2^{k+1}-1]$ (1/128--64\,Hz). With
$\ell_k=\sum_r\mathrm{num}_{k,r}/\sum_r\mathrm{den}_r$, $w_k=\max(\ell_k,0)$ and $\nu_k=2^{k+0.5}/128$\,Hz,
$\nu_c=\exp\bigl(\sum_kw_k\ln\nu_k/\sum_kw_k\bigr)$. It is missing if $\sqrt{\max(\sum_k\ell_k,0)}<0.05$. Unlike the RXTE
version (\cref{eq:nuc}), the weights are not divided by the bin width; for octave bins this matters only for the lowest bin.
\item \textbf{State.} \Cref{eq:state} is applied to the 0.1--10\,Hz rms, the full band of \citet{RemillardMcClintock2006}.
\item \textbf{Status.} \code{ok} needs at least 3 valid segments and \code{rate_2_10}~$\ge20$\,c\,s$^{-1}$; otherwise
\code{missing} or \code{excluded} (\cref{tab:qualityrules}).
\end{itemize}

\paragraph{Background proxy (v10, \code{v10lib.bkg_ratio}).} $\beta=N(\mathrm{PI}\,1200\text{--}1499)/N(\mathrm{PI}\,200\text{--}999)$
over the whole file, the ratio of 12--15\,keV to 2--10\,keV events. Above 12\,keV the source contributes little, so a high
$\beta$ indicates background. The threshold, frozen at 0.2687, is the 99th percentile of $\beta$ among observations brighter
than 500\,c\,s$^{-1}$, where the background is negligible. It was used only for sensitivity analyses in v10 and v11.

\paragraph{Per-bin errors and Lorentzian fits (v11c, \code{v11lib}).}\label{sec:lorentz} \code{lb_se} gives a delta-method
error for each octave-bin variance $\ell_k$. \code{fit_lorentz} fits one or two zero-centred Lorentzians whose integrated
variance in bin $[\nu^\mathrm{lo}_k,\nu^\mathrm{hi}_k]$ (edges $(j\mp0.5)/128$\,Hz) is
\begin{equation}
L_k(r^2,\Delta)=r^2\,\frac{2}{\pi}\Bigl[\arctan\frac{\nu^\mathrm{hi}_k}{\Delta}-\arctan\frac{\nu^\mathrm{lo}_k}{\Delta}\Bigr],\qquad\Delta\in[1/128,\,128]\ \mathrm{Hz}.\label{eq:lorentz}
\end{equation}
The fit is weighted least squares in log-parameters (\code{scipy.optimize.least_squares}), with four starting points per
model; at least 6 usable bins are required. With $\mathrm{BIC}=\chi^2+k\ln n$, two components are chosen if
$\mathrm{BIC}_1-\mathrm{BIC}_2>6$. The width of the lowest-frequency component is $\nu_L$.

\paragraph{Reading modes and storage.} v9 parsed only a prefix of each file, stopping at 6 segments (\cref{tab:nicermodes}).
v10 appended the tail up to 1\,GB and re-parsed prefix and tail as one stream (\code{v10lib.read_full}); 366 of the 368 N1
observations remained \code{ok}. v11 and v13 streamed whole files into the frozen \code{obs_features} without storing them
(\code{v11lib.obs_features_stream}), with 64\,MB requests in windows of 6 concurrent requests consumed in order. The fetch
scripts keep a checkpoint (\file{observations_partial.jsonl}), so an interrupted run resumes without re-reading finished
observations.

% -----------------------------------------------------------------------------------------------------
\subsection{NICER background: 3C50 correction and screening (v12, v13)}\label{sec:3c50}
\Cref{fig:3c50flow} shows the background step. It is applied only to observations fainter than 500\,c\,s$^{-1}$ in
2--10\,keV; for brighter ones the background is treated as negligible ($f=0$).

\begin{figure}[htbp]
\centering
\begin{tikzpicture}[x=1cm, y=1cm,
  st/.style={fc proc, text width=2.62cm, minimum height=16mm},
  badge/.style={fc ver, anchor=south west}]
  \node[fc data, text width=2.62cm, minimum height=16mm] (in) at (1.45,0) {accepted NICER obs.\ (\code{ok}) with rate $<500$\,c\,s$^{-1}$};
  \node[st] (dl) at (4.62,0) {WSL: \code{curl} \code{_cl} and \code{_ufa} files (3 tries); SHA256 to the log};
  \node[st] (bg) at (7.79,0) {\code{nibackgen3C50}, CALDB library, \code{gainepoch=AUTO}};
  \node[st] (sp) at (10.96,0) {keep \file{tot_<o>.pi}, \file{bkg_<o>.pi}; delete event files};
  \node[st] (fr) at (14.13,0) {band fractions $f_\mathcal{X}=\sum B/\sum T$ for $\mathcal X$ = A, B, C, T};
  \foreach \a/\b in {in/dl,dl/bg,bg/sp,sp/fr} \draw[fc arr] (\a) -- (\b);
  \node[fc dec, text width=1.9cm, aspect=1.9] (q) at (14.13,-2.55) {3C50 ok and $f_\mathrm{T}\le0.20$?};
  \node[st] (co) at (10.2,-2.55) {correct rms by $1/(1-f_\mathrm{T})$, colours band by band (Eq.~\ref{eq:3c50}); new state};
  \node[fc no, text width=2.62cm, minimum height=10mm] (ex) at (14.13,-4.5) {excluded: 3C50 failed or $f_\mathrm{T}>0.20$};
  \node[fc, draw=inkgrey, fill=boxgrey, text width=3.1cm, minimum height=16mm] (out) at (6.3,-2.55) {\file{nicer_corrected_screened.csv} (v12); v13 test table};
  \node[fc data, text width=2.62cm, minimum height=12mm] (br) at (1.45,-2.55) {bright obs.: $f=0$, kept unchanged};
  \draw[fc arr] (fr) -- (q); \draw[fc arr] (q) -- node[fc note, above] {yes} (co); \draw[fc arr] (q) -- node[fc note, right] {no} (ex);
  \draw[fc arr] (co) -- (out); \draw[fc arr] (br) -- (out);
  \foreach \n/\v in {in/v12, dl/v12, bg/v12, sp/v12, fr/v12, co/{v12, frozen in v13}}
     \node[badge] at ([xshift=2pt, yshift=-0.4pt]\n.north west) {\v};
\end{tikzpicture}
\caption{3C50 background step (\file{58_v12_bkg.py}, \file{62_v13_bkg.py}, \file{v12_bkg3c50.sh}; correction in
\code{corrected} of \file{59_v12_analysis.py}, copied unchanged into \file{63_v13_analysis.py}).}\label{fig:3c50flow}
\end{figure}

\paragraph{Running 3C50.} The 3C50 model \citep{Remillard2022} estimates the background spectrum of an observation from a
library of blank-sky observations, selected by background-sensitive count rates measured in the observation itself (hence
the unfiltered \code{_ufa} file). For each faint observation the Python driver calls, through \code{wsl -d Ubuntu-22.04}, the shell
script \file{v12_bkg3c50.sh} with the ObsID, the cleaned and unfiltered URLs, an output folder and the CALDB folder. The
script activates the HEASoft conda environment and points \code{CALDB}, \code{CALDBCONFIG} and \code{CALDBALIAS} at the
unpacked library. It downloads the two files into \file{/tmp/bkg3c50/<obsid>/xti/event_cl/} and unzips them. It then runs
\begin{verbatim}
nibackgen3C50 rootdir=/tmp/bkg3c50 obsid=<obsid>
              bkgidxdir=CALDB bkglibdir=CALDB gainepoch=AUTO clobber=yes chatter=1
              totspec=tot_<obsid> bkgspec=bkg_<obsid>
\end{verbatim}
copies the two spectra out and deletes the working folder. The driver reads the spectra (\code{RATE}, or \code{COUNTS}
divided by \code{EXPOSURE}) and sums them in the PI bands A, B, C and T (200--999). It records the band rates $T_\mathcal X$
and $B_\mathcal X$, the fractions $f_\mathcal X=B_\mathcal X/T_\mathcal X$ and the 3C50 exposure in
\file{data/v12/processed/bkg3c50.csv}. Six workers run in parallel, and a checkpoint file makes the step resumable.

\paragraph{Failures.} In v12, 362 of the 503 faint observations gave spectra; 141 failed, mostly because the tool rejects
data processed with the oldest gain calibration. In v13, 109 of 148 succeeded. Failed observations are excluded from the
corrected analyses.

\paragraph{Correction and screening.} An observation is kept if it is bright, or if 3C50 succeeded with
$f_\mathrm{T}\le0.20$; thresholds 0.10 and 0.30 are sensitivity variants. For kept faint observations,
\begin{equation}
T_k'=\frac{T_k}{1-f_\mathrm{T}}\quad(k=1,2,3,\text{state}),\qquad c_1'=\frac{B(1-f_B)}{A(1-f_A)},\qquad c_2'=\frac{C(1-f_C)}{B(1-f_B)},\label{eq:3c50}
\end{equation}
and the state is re-classified with the corrected state rms. A fractional rms is normalised to the total mean count rate;
renormalising it to the source mean $S=(1-f)(S+B)$ multiplies it by $1/(1-f)$, assuming a constant background. $\nu_c$ is
not changed, because it depends only on the shape of the power spectrum. A check requires the 3C50 total 2--10\,keV rate to
be within 0.7--1.3 of the pipeline rate for at least 90\% of observations.

\paragraph{GTI-restricted variant (v13, descriptive).} For the faint v13 observations, the cleaned event files were kept
(\code{V13_KEEP_CL}). Their features were recomputed using only events inside the GTIs of the 3C50 total spectrum
(\code{gti_restricted} in \file{63_v13_analysis.py}), to check that the two time selections give the same picture.

% -----------------------------------------------------------------------------------------------------
\subsection{MAXI/GSC: daily light curves (v7c, v8d, v9c)}\label{sec:maxi}
MAXI points are days, not pointed observations. \Cref{fig:maxiflow} shows the chain.

\begin{figure}[htbp]
\centering
% MAXI/GSC pipeline (v7c, v8d, v9c) and the version that added each step. Copied from report_en/figures/flowcharts (styles in workflow_report.tex).
\begin{tikzpicture}[x=1cm, y=1cm,
  st/.style={fc proc, text width=2.62cm, minimum height=15mm},
  badge/.style={fc ver, anchor=south west}]
  % row 1: strict pipeline (v7c2)
  \node[fc data, text width=2.62cm, minimum height=15mm] (lab) at (1.45,0) {labels frozen first: confirmed BHs, non-pulsing NSs, pulsars};
  \node[st] (mt)  at (4.62,0)  {MAXI sources matched by Sesame position ($\le0.1^\circ$); blends excluded};
  \node[fc data, text width=2.62cm, minimum height=15mm] (lc) at (7.79,0) {1-day light curves $L$, $M$, $H$ (2--4, 4--10, 10--20\,keV)};
  \node[st] (cut) at (10.96,0) {cuts ($\ge\sqrt3\,\sigma$, outliers, $\ge100$ points): 13 BH, 48 NPNS, 39 PSR};
  \node[st] (ft)  at (14.13,0) {SC, HC, RelInt, $C_4$; sets CCI, C2D, C1D (Eq.~\ref{eq:maxi})};
  \foreach \a/\b in {lab/mt,mt/lc,lc/cut,cut/ft} \draw[fc arr] (\a) -- (\b);
  % row 2: absorption correction (v8d) feeding the same models
  \node[fc data, text width=2.62cm, minimum height=15mm] (nh) at (4.62,-2.35) {HI4PI Galactic $N_\mathrm{H}$ at each source};
  \node[st] (tr)  at (7.79,-2.35) {\texttt{tbabs}$\times$power-law transmission $t_L,t_M,t_H$};
  \node[st] (ftc) at (10.96,-2.35) {corrected fluxes $L/t_L$, \dots: C2D$_{N_\mathrm{H}}$, C1D$_{N_\mathrm{H}}$};
  \node[fc model, text width=2.62cm, minimum height=15mm] (md) at (14.13,-2.35) {LR, RF, kNN, SVM; LOSO over sources; $\le200$ training points per source};
  \draw[fc arr] (nh) -- (tr); \draw[fc arr] (tr) -- (ftc); \draw[fc arr] (ftc) -- (md);
  \draw[fc arr] (ft) -- (md);
  \draw[fc arr] (lc.south) -- ++(0,-0.5) -| ([xshift=-6mm]ftc.north);
  % row 3: strata (v9c) and the reproduction (v7c1)
  \node[st, text width=2.62cm, minimum height=11mm] (ter) at (14.13,-4.45) {points split into tertiles of $C_4$; gain within each};
  \draw[fc arr] (md) -- (ter);
  \node[fc data, text width=2.62cm, minimum height=11mm] (db) at (1.45,-4.45) {de Beurs et al.\ (2022) processed colours, 44 systems};
  \node[fc model, text width=2.62cm, minimum height=11mm] (dbm) at (4.62,-4.45) {their SVM and kNN settings, 10 subsamples};
  \node[fc gate, text width=2.62cm, minimum height=11mm] (dbr) at (7.79,-4.45) {reproduction within the pre-set tolerance};
  \draw[fc arr] (db) -- (dbm); \draw[fc arr] (dbm) -- (dbr);
  % badges
  \foreach \n/\v in {lab/v7c2, mt/v7c2, lc/v7c2, cut/v7c2, ft/v7c2, md/v7c2, nh/v8d, tr/v8d, ftc/v8d, ter/v9c, db/v7c1, dbm/v7c1, dbr/v7c1}
     \node[badge] at ([xshift=2pt, yshift=-0.4pt]\n.north west) {\v};
\end{tikzpicture}

\caption{MAXI/GSC pipeline: the strict pipeline (v7c2), the absorption correction (v8d), the colour strata (v9c) and the
reproduction of \citet{deBeurs2022} (v7c1), which uses their processed colours instead of this pipeline.}\label{fig:maxiflow}
\end{figure}

\paragraph{Source matching (\file{38_v7c2_select.py}).}
\begin{enumerate}
\item The label lists are frozen first: 25 BHs \citep{Marcel2026}, the 115 MINBAR bursters and the pulsar lists. Sources
with a pulse period in the Liu et al.\ HMXB or LMXB catalogues (CDS J/A+A/455/1165 and J/A+A/469/807) or in Patruno \&
Watts form the separate pulsar class. MINBAR bursters without a pulse period are the ``non-pulsing'' NSs (NPNS).
\item The MAXI source index is scraped from the 13 category pages \file{http://maxi.riken.jp/top/lc_<cat>.html}. Each entry
gives a name and a J-identifier, whose digits encode an approximate position (RA hhmm, Dec $\pm$dd.d).
\item Entries whose name contains ``with'' are blends of two sources and are removed. The remaining MAXI names are resolved
with Sesame. A label source matches a MAXI source within $0.1^\circ$ of the Sesame position, or within $0.3^\circ$ of the
J-name position if Sesame fails.
\item A MAXI source matched by label sources of different classes, or by several label sources that are not aliases of one
object, is excluded as ambiguous (\file{data/v7c/c2_sources_ambiguous.csv}).
\end{enumerate}

\paragraph{Light curves and cuts.} For every matched source the 1-day light curve
\file{<J>_g_lc_1day_all.dat} is downloaded (columns MJD, total flux and error, then $L$, $M$, $H$ with errors for 2--4,
4--10 and 10--20\,keV). The cuts of \citet{deBeurs2022} are then applied:
\begin{itemize}
\item every band must have flux/error $\ge\sqrt3$ and a positive error;
\item points with $L+M+H\ge$ mean $+\,10$\,sd of the source are dropped;
\item the source needs at least 100 remaining points.
\end{itemize}
Result: \file{data/v7c/maxi_points.csv} with 132\,057 points of 13 BH, 48 NPNS and 39 pulsar sources.

\paragraph{Features (\file{39_v7c2_analysis.py}).} Per point,
\begin{equation}
\mathrm{SC}=\frac{M-L}{M+L},\qquad\mathrm{HC}=\frac{H-L}{H+L},\qquad C_4=\frac{H-M}{H+M},\qquad
\mathrm{RelInt}=\frac{L+M+H}{Q_{99.99}(L+M+H)},\label{eq:maxi}
\end{equation}
where $Q_{99.99}$ is the 99.99th percentile over that source's points. The feature sets are CCI = (SC, HC, RelInt),
C2D = (SC, HC) and C1D = ($C_4$), the only colour without the 2--4\,keV band.

\paragraph{Absorption correction (v8d, \file{45_v8d_nh.py}, \file{v8d_heasoft.sh}).}
\begin{enumerate}
\item HEASoft \code{nh} gives the HI4PI column at each source position (\code{disio=0.1}; the weighted average
\code{avwnh} is used).
\item XSPEC computes photon fluxes of \code{tbabs*powerlaw} ($\Gamma=2$; \code{abund wilm}, \code{xsect vern}; a dummy
response from 0.5 to 30\,keV) in the three MAXI bands. The grid is $N_\mathrm{H}=0$ plus 200 logarithmic steps from
$10^{19}$ to $10^{23.5}$\,cm$^{-2}$.
\item The transmission is $t_\mathcal X(N_\mathrm H)=\Phi_\mathcal X(N_\mathrm H)/\Phi_\mathcal X(0)$, interpolated linearly
in $\log N_\mathrm H$.
\item The corrected fluxes $L/t_L$, $M/t_M$, $H/t_H$ give SC$'$, HC$'$, $C_4'$ and RelInt$'$ through \cref{eq:maxi} (sets
C2D$_{N_\mathrm H}$, C1D$_{N_\mathrm H}$, CCI$_{N_\mathrm H}$).
\end{enumerate}

\paragraph{Colour strata (v9c, \file{50_v9c_maxi_strata.py}).} The points are split into tertiles of $C_4$, and the
models are compared within each tertile.

\paragraph{Reproduction of de Beurs et al.\ (v7c1, \file{37_v7c1_repro.py}).} Their processed colours of 44 systems and
their ten 20\% subsamples are used with their three-class settings (BH, NPNS, pulsar): kNN with $k=24$ and SVM with $C=0.655$,
$\gamma=0.585$. The reproduction counts as successful if, for each algorithm, the number of correctly classified systems
is within 2 of the published number (kNN 37, SVM 36) and the number of correct BHs within 1 (6 and 5)
(\code{V7_REPRO_TOL_SOURCES}, \code{V7_REPRO_TOL_BH}). Otherwise the script stops.
